\begin{resume}
    \section*{1.\hspace{0.75em}基本情况}
    \anonrvwinfo{周广义}{XXX}，男，山东菏泽人，2001年10月出生，西安电子科技大学\anonrvwinfo{计算机科学与技术}{XXX}学院\anonrvwinfo{软件工程}{XXX}专业2023级硕士研究生。
    \section*{2.\hspace{0.75em}教育背景}
    \begin{resumelist*}
        \resumelistitem 2019.09～2023.06，长安大学，本科，专业：\anonrvwinfo{计算机科学与技术}{XXX}
        \resumelistitem 2023.09～\phantom{2025.07}，西安电子科技大学，硕士研究生，专业：\anonrvwinfo{软件工程}{XXX}
    \end{resumelist*}
    \section*{3.\hspace{0.75em}攻读硕士学位期间的研究成果}
    \begin{resumelist}{\hspace{-0.25em}3.1\hspace{0.5em} 发表专利}
        \resumelistitem
        王建东,周广义,张涛,王绍青,黄思宁,潘晟冉,刘永昌,卢天增,李承浩,毕冒泽. 一种多维感知注视点渲染方法、系统和介质：中国, 2026104215486[P]. 2026-04-01.
        \resumelistitem
        王建东,潘晟冉,张琛,张志为,黄思宁,周广义,李承浩,韩雨. 一种基于多维特征的自适应混合渲染方法、系统及介质：中国, 2026102999624[P]. 2026-03-12.
        \resumelistitem
        徐扬,黄思宁,张志为,潘晟冉,薛天琳,周广义,李承浩,亢玉锋,刘庆元.一种面向复杂语义的渐进式文本引导三维模型生成方法及系统：中国, 2026102999554[P]. 2026-03-12.
    \end{resumelist}
    \begin{resumelist}{\hspace{-0.25em}3.2\hspace{0.5em} 软件著作权}
        \resumelistitem
        西安电子科技大学青岛计算技术研究院，周广义，黄思宁，潘晟冉.智慧楼宇孪生系统 V1.0. 2025SR0379713.
    \end{resumelist}
    \begin{resumelist}{\hspace{-0.25em}3.3\hspace{0.5em} 参与科研项目及获奖} 
        %		\resumelistitem XXX项目, 项目名称, 起止时间, 完成情况, 作者贡献。
        %\anonrvwinfo{XXX, XXX, XXX等}{第一作者}.
        \resumelistitem 委托开发项目，智慧青研院管理可视化平台，2024.07-2024.09，结项，主要研发人员。

        \resumelistitem 企业委托项目，智能工厂生产线监控平台，2024.09-2024.12，结项，主要研发人员。

        \resumelistitem 企业委托项目，金圃广源数字工厂管理系统，2025.01-2025.5，完成，主要研发人员。

        \resumelistitem 委托开发项目，三维可视化低代码平台，2024.09-至今，在研，主要研发人员。

    \end{resumelist}
\end{resume}